% Options for packages loaded elsewhere
% Options for packages loaded elsewhere
\PassOptionsToPackage{unicode}{hyperref}
\PassOptionsToPackage{hyphens}{url}
\PassOptionsToPackage{dvipsnames,svgnames,x11names}{xcolor}
%
\documentclass[
  french,
  11pt,
  letterpaper,
  DIV=11,
  numbers=noendperiod]{scrartcl}
\usepackage{xcolor}
\usepackage[top=30mm,left=25mm,right=25mm,bottom=30mm]{geometry}
\usepackage{amsmath,amssymb}
\setcounter{secnumdepth}{-\maxdimen} % remove section numbering
\usepackage{iftex}
\ifPDFTeX
  \usepackage[T1]{fontenc}
  \usepackage[utf8]{inputenc}
  \usepackage{textcomp} % provide euro and other symbols
\else % if luatex or xetex
  \usepackage{unicode-math} % this also loads fontspec
  \defaultfontfeatures{Scale=MatchLowercase}
  \defaultfontfeatures[\rmfamily]{Ligatures=TeX,Scale=1}
\fi
\usepackage{lmodern}
\ifPDFTeX\else
  % xetex/luatex font selection
\fi
% Use upquote if available, for straight quotes in verbatim environments
\IfFileExists{upquote.sty}{\usepackage{upquote}}{}
\IfFileExists{microtype.sty}{% use microtype if available
  \usepackage[]{microtype}
  \UseMicrotypeSet[protrusion]{basicmath} % disable protrusion for tt fonts
}{}
\makeatletter
\@ifundefined{KOMAClassName}{% if non-KOMA class
  \IfFileExists{parskip.sty}{%
    \usepackage{parskip}
  }{% else
    \setlength{\parindent}{0pt}
    \setlength{\parskip}{6pt plus 2pt minus 1pt}}
}{% if KOMA class
  \KOMAoptions{parskip=half}}
\makeatother
% Make \paragraph and \subparagraph free-standing
\makeatletter
\ifx\paragraph\undefined\else
  \let\oldparagraph\paragraph
  \renewcommand{\paragraph}{
    \@ifstar
      \xxxParagraphStar
      \xxxParagraphNoStar
  }
  \newcommand{\xxxParagraphStar}[1]{\oldparagraph*{#1}\mbox{}}
  \newcommand{\xxxParagraphNoStar}[1]{\oldparagraph{#1}\mbox{}}
\fi
\ifx\subparagraph\undefined\else
  \let\oldsubparagraph\subparagraph
  \renewcommand{\subparagraph}{
    \@ifstar
      \xxxSubParagraphStar
      \xxxSubParagraphNoStar
  }
  \newcommand{\xxxSubParagraphStar}[1]{\oldsubparagraph*{#1}\mbox{}}
  \newcommand{\xxxSubParagraphNoStar}[1]{\oldsubparagraph{#1}\mbox{}}
\fi
\makeatother


\usepackage{longtable,booktabs,array}
\newcounter{none} % for unnumbered tables
\usepackage{calc} % for calculating minipage widths
% Correct order of tables after \paragraph or \subparagraph
\usepackage{etoolbox}
\makeatletter
\patchcmd\longtable{\par}{\if@noskipsec\mbox{}\fi\par}{}{}
\makeatother
% Allow footnotes in longtable head/foot
\IfFileExists{footnotehyper.sty}{\usepackage{footnotehyper}}{\usepackage{footnote}}
\makesavenoteenv{longtable}
\usepackage{graphicx}
\makeatletter
\newsavebox\pandoc@box
\newcommand*\pandocbounded[1]{% scales image to fit in text height/width
  \sbox\pandoc@box{#1}%
  \Gscale@div\@tempa{\textheight}{\dimexpr\ht\pandoc@box+\dp\pandoc@box\relax}%
  \Gscale@div\@tempb{\linewidth}{\wd\pandoc@box}%
  \ifdim\@tempb\p@<\@tempa\p@\let\@tempa\@tempb\fi% select the smaller of both
  \ifdim\@tempa\p@<\p@\scalebox{\@tempa}{\usebox\pandoc@box}%
  \else\usebox{\pandoc@box}%
  \fi%
}
% Set default figure placement to htbp
\def\fps@figure{htbp}
\makeatother



\ifLuaTeX
\usepackage[bidi=basic,shorthands=off]{babel}
\else
\usepackage[bidi=default,shorthands=off]{babel}
\fi
\ifLuaTeX
  \usepackage{selnolig} % disable illegal ligatures
\fi


\setlength{\emergencystretch}{3em} % prevent overfull lines

\providecommand{\tightlist}{%
  \setlength{\itemsep}{0pt}\setlength{\parskip}{0pt}}



 


\usepackage{graphicx}
\KOMAoption{captions}{tableheading}
\makeatletter
\@ifpackageloaded{caption}{}{\usepackage{caption}}
\AtBeginDocument{%
\ifdefined\contentsname
  \renewcommand*\contentsname{Table des matières}
\else
  \newcommand\contentsname{Table des matières}
\fi
\ifdefined\listfigurename
  \renewcommand*\listfigurename{Liste des Figures}
\else
  \newcommand\listfigurename{Liste des Figures}
\fi
\ifdefined\listtablename
  \renewcommand*\listtablename{Liste des Tables}
\else
  \newcommand\listtablename{Liste des Tables}
\fi
\ifdefined\figurename
  \renewcommand*\figurename{Figure}
\else
  \newcommand\figurename{Figure}
\fi
\ifdefined\tablename
  \renewcommand*\tablename{Table}
\else
  \newcommand\tablename{Table}
\fi
}
\@ifpackageloaded{float}{}{\usepackage{float}}
\floatstyle{ruled}
\@ifundefined{c@chapter}{\newfloat{codelisting}{h}{lop}}{\newfloat{codelisting}{h}{lop}[chapter]}
\floatname{codelisting}{Listing}
\newcommand*\listoflistings{\listof{codelisting}{Liste des Listings}}
\makeatother
\makeatletter
\makeatother
\makeatletter
\@ifpackageloaded{caption}{}{\usepackage{caption}}
\@ifpackageloaded{subcaption}{}{\usepackage{subcaption}}
\makeatother
\usepackage{bookmark}
\IfFileExists{xurl.sty}{\usepackage{xurl}}{} % add URL line breaks if available
\urlstyle{same}
% fallback for those not using the hyperref driver hyperxmp:
\makeatletter
\@ifundefined{xmpquote}{\newcommand{\xmpquote}[1]{#1}}{}
\makeatother
\hypersetup{
  pdflang={fr},
  colorlinks=true,
  linkcolor={blue},
  filecolor={Maroon},
  citecolor={Blue},
  urlcolor={Blue},
  pdfcreator={LaTeX via pandoc}}


\author{}
\date{}
\begin{document}


\begin{titlepage}
\centering


% Logo
\includegraphics[width=7cm]{ensai-logo.png}\par

\vspace{1.5cm}

% Titre principal
{\Huge Rapport final - Projet informatique \par}
\vspace{0.5cm}
{\Large 2A \par}

\vspace{1.5cm}

% Sous-titre / Nom de l'application
{\LARGE \textbf{Concurrent SNCF} \par}

\vspace{1cm}

% Bloc Étudiants (aligné à gauche)
\begin{flushleft}
\large
\textbf{Étudiants :}\\
Adam BENAMER\\
Satya KREMER\\
Nathan TIATIA\\
Laurine LEROY\\
Lou-Ann CLEMENT\\
Nour NSIRI
\end{flushleft}

\vspace{0.05cm}

% Bloc Encadrant 
\begin{flushright}
\large
\textbf{Encadrant pédagogique :}\\
KEVIN LEROY
\end{flushright}


\vfill

% Année universitaire
{\normalsize Année universitaire 2026--2027 \par}

\end{titlepage}
\tableofcontents
\newpage

\section{Introduction}\label{introduction}

Dans un monde exigeant une grande mobilité des personnes et des biens,
se déplacer est devenu un enjeu majeur pour tous et toutes. Cependant,
ces déplacements viennent souvent à l'encontre de préoccupations
environnementales grandissantes. En effet, les trajets en avion ou même
en voiture sont extrêmement polluants : il est donc pertinent de se
pencher sur une de leurs alternatives la plus verte, c'est-à-dire le
train. En effet, le train est une industrie inestimable pour faire face
aux enjeux climatiques de notre monde. Ainsi, il est tout aussi
nécessaire d'avoir des services qui font fonctionner la concurrence,
afin de permettre aux utilisateurs des chemins de fer d'avoir la
meilleure expérience possible. C'est pour cette raison que la création
d'un concurrent à la SNCF s'impose comme un sujet particulièrement
important, et demandant un traitement attentif et complet.

Afin d'accomplir cela, nous avons développé une application qui vise à
mettre en place le plus grand nombre possible de fonctionnalités
proposées par la SNCF : ce rapport vise donc à expliquer dans le détail
comment cette application a été conçue, structurée et codée par notre
équipe. Nous verrons dans un premier temps le contexte de l'application
et le cahier des charges dressé, ainsi que l'organisation de travail
mise en place au sein de l'équipe. Puis nous nous pencherons sur la
modélisation de notre application et son architecture générale. Nous
nous intéresserons ensuite à l'implémentation, notamment par le biais de
diagramme de classes, afin d'expliciter les différentes couches de code
mises en place et leurs interactions. Nous verronns aussi à cette
occasion les outils utilisés et les tests mis en place dans une démarche
d'assurance-qualité. Enfin, lors d'une dernière partie de prise de recul
nous verrons les possibles modifications à notre application selon
différents scénarios.

\section{I. Présentation du projet}\label{i.-pruxe9sentation-du-projet}

\subsection{Cahier des charges et fonctionnalités de
l'application}\label{cahier-des-charges-et-fonctionnalituxe9s-de-lapplication}

Au cours de ce semestre, nous avons donc développé une application qui
cherche à approcher au mieux les fonctionnalités proposées par la SNCF.
Celles-ci sont constituées de cinq fonctionnalités de base et deux
fonctionnalités optionnelles. Ces fonctionnalités utilisent parmi les
données transmises par l'API SNCF les gares, lignes, trajets et prix
présents dans la base de données. La récupération, l'exploitation et la
transmission à notre base de données de ces données SNCF seront
explicités lors de l'explication des couches de code.

La première fonctionnalité est une fonctionnalité de gestion des comptes
utilisateurs, afin d'avoir les différents niveaux d'habilitation
nécessaires à l'application. Ceux-ci sont au nombre de trois : les
administrateurs sont les seuls à pouvoir changer le niveau
d'habilitation d'autres comptes, mais ont accès à toutes les autres
fonctionnalités accessibles aux deux autres catégories de compte, les
collaborateurs et les clients. Tandis que les collaborateurs peuvent
créer des lignes et trajets, les clients ne peuvent que rechercher ces
lignes et trajets, et réserver les trajets disponibles.

La deuxième fonctionnalité demandée est celle de gestion des lignes
d'exploitation. Cela est effectué par les comptes collaborateurs, en
utilisant l'API externe de la SNCF. Pour faire cela, les collaborateurs
recherchent d'abord des gares existantes dans la base de données de la
SNCF. Cette recherche permet ensuite au collaborateur d'utiliser les
gares obtenues pour vérifier l'existence d'une ligne dans la base de
données de la SNCF, et de créer ladite ligne.

La troisième fonctionnalité que nous devons implémenter est celle de
gestion des trajets proposés. Un trajet se trouve sur une ligne, et
possède donc une date et une heure de départ, un nombre de places et un
tarif. Le temps de trajet est calculé avec l'API SNCF, tandis que le
nombre de places est une variable à ajuster par le collaborateur créant
le trajet. En plus de cela, il est important de vérifier qu'un trajet
peut bel et bien être effectué entre deux gares, et qu'il existe donc
une ligne déjà créée par le compte collaborateur qui permet ce trajet.
Enfin, les comptes collaborateurs sont capables de modifier ou bien de
supprimer un trajet.

La quatrième fonctionnalité est celle de recherche de trajet : un compte
client peut rechercher les trajets créés grâce à la fonctionnalité
précédente, en entrant une gare d'arrivée et une gare de départ. A
l'issue de cette recherche, toutes les informations nécessaires lors de
la création d'un trajet sont visibles par le client, à l'exception du
tarif et des places (qui peuvent être néanmoins visibles lors de la
réservation d'un trajet).

La cinquième et dernière fonctionnalité est une fonctionnalité de
positionnement commercial, laissée libre. Nous avons choisi
d'implémenter la tarification transparente : ainsi, lors de la
réservation d'un trajet, un client peut voir les différentes composantes
du prix total. Nous avons aussi implémenté lors de la réservation la
visibilité du nombre de places restantes dans la classe voulue (première
ou deuxième).

De plus, nous avons aussi développé les deux fonctionnalités
optionnelles proposées par le sujet : celle des classes de voyage, et la
fonctionnalité de réservation. Cette dernière permet à un compte client
de réserver et annuler une réservation, ainsi que de visualiser ses
réservations.

\subsection{Organisation de l'équipe}\label{organisation-de-luxe9quipe}

L'organisation du groupe a été prise en charge par Lou-Ann Clement en
tant que cheffe de projet, tandis que les questions purement techniques
portant sur le code étaient gérées par Satya Kremer, le tech leader du
groupe.

Les tâches étaient réparties lors de réunions de travail : lors de la
phase d'analyse les diagrammes et parties du rapport ont été attribuées
lors de ces réunions, et lors de la partie code plusieurs réunions ont
euu lieu pour répartir les tâches. Au début celle-ci se présentaient
sous la forme de classes à coder à travers plusieurs couches, et au fur
et à mesure de l'avancement du projet ces tâches étaient des erreurs ou
corrections remarquées par des membres de l'équipe, et distribuées à
tous selon les capacités techniques de chacun. Pour améliorer la
communication et la collaboration du groupe en dehors de ces réunions,
un groupe Whatsapp a été utilisé régulièrement, ainsi qu'un Google Drive
pour centraliser diagrammes et retours du tuteur notamment.;

\section{II. Modélisation}\label{ii.-moduxe9lisation}

\section{III. Implémentation}\label{iii.-impluxe9mentation}

\subsection{Outils utilisés}\label{outils-utilisuxe9s}




\end{document}
